\documentclass[aspectratio=169]{beamer}

% ------------------------------------------------
% Pacotes
% ------------------------------------------------
\usepackage[utf8]{inputenc}
\usepackage[T1]{fontenc}
\usepackage[brazil]{babel}
\usepackage{lmodern}

\usepackage{amsmath}
\usepackage{amssymb}
\usepackage{graphicx}
\usepackage{booktabs}
\usepackage{multicol}
\usepackage{xcolor}
\usepackage{hyperref}

% ------------------------------------------------
% Tema
% ------------------------------------------------
\usetheme{Madrid}
\usecolortheme{default}

% Remove ícones de navegação
\setbeamertemplate{navigation symbols}{}

% ------------------------------------------------
% Rodapé
% ------------------------------------------------
\setbeamertemplate{footline}{
    \leavevmode%
    \hbox{%
        \begin{beamercolorbox}[
            wd=.80\paperwidth,
            ht=2.5ex,
            dp=1ex,
            leftskip=1em
        ]{author in head/foot}
            \insertshorttitle
        \end{beamercolorbox}%
        \begin{beamercolorbox}[
            wd=.20\paperwidth,
            ht=2.5ex,
            dp=1ex,
            center
        ]{date in head/foot}
            \insertframenumber/\inserttotalframenumber
        \end{beamercolorbox}%
    }
}

% ------------------------------------------------
% Cores
% ------------------------------------------------
\definecolor{academicblue}{RGB}{25,65,110}
\definecolor{lightgray}{RGB}{245,245,245}

\setbeamercolor{structure}{
    fg=academicblue
}

\setbeamercolor{frametitle}{
    fg=white,
    bg=academicblue
}

% Capa
\setbeamercolor{title}{
    fg=white,
    bg=academicblue
}

\setbeamercolor{subtitle}{
    fg=white,
    bg=academicblue
}

\setbeamercolor{author}{
    fg=black
}

\setbeamercolor{institute}{
    fg=black
}

\setbeamercolor{date}{
    fg=black
}

% Blocos
\setbeamercolor{block title}{
    fg=white,
    bg=academicblue
}

\setbeamercolor{block body}{
    fg=black,
    bg=lightgray
}

% ------------------------------------------------
% Informações da apresentação
% ------------------------------------------------
\title[aaasadffsaa]{
testeaaaa
}

\subtitle{Subtítulo, se necessário}

\author{Nome do Autor - Matrícula}

\institute{
    Universidade / Instituição \\
    Curso ou Programa \\
    Grupo de Pesquisa
}

\date{2026}

% Para adicionar um logo:
%
% \titlegraphic{
%     \includegraphics[width=2cm]{logo.png}
% }

% ------------------------------------------------
% Documento
% ------------------------------------------------
\begin{document}

% =================================================
% CAPA
% =================================================
\begin{frame}
    \titlepage
\end{frame}

% =================================================
% SUMÁRIO
% =================================================
\begin{frame}{Sumário}
    \tableofcontents
\end{frame}

% =================================================
% INTRODUÇÃO
% =================================================
\section{Introdução}

\begin{frame}{Contextualização}

\begin{itemize}
    \item Apresente brevemente o problema estudado.
    \item Explique o contexto no qual o problema ocorre.
    \item Mostre por que esse problema é relevante.
    \item Apresente a principal lacuna identificada.
\end{itemize}

\vspace{0.4cm}

\begin{block}{Questão de pesquisa}
Qual é a principal pergunta que este trabalho busca responder?
\end{block}

\end{frame}

% =================================================
\begin{frame}{Motivação}

\begin{columns}[T]

    \begin{column}{0.48\textwidth}

        \textbf{Problema}

        \begin{itemize}
            \item Problema 1
            \item Problema 2
            \item Problema 3
        \end{itemize}

    \end{column}

    \begin{column}{0.48\textwidth}

        \textbf{Impactos}

        \begin{itemize}
            \item Impacto 1
            \item Impacto 2
            \item Impacto 3
        \end{itemize}

    \end{column}

\end{columns}

\end{frame}

% =================================================
% OBJETIVOS
% =================================================
\section{Objetivos}

\begin{frame}{Objetivos}

\begin{block}{Objetivo geral}
Descrever aqui o objetivo principal do trabalho.
\end{block}

\vspace{0.3cm}

\textbf{Objetivos específicos}

\begin{itemize}
    \item Investigar ...
    \item Modelar ...
    \item Desenvolver ...
    \item Comparar ...
    \item Avaliar ...
\end{itemize}

\end{frame}

% =================================================
% METODOLOGIA
% =================================================
\section{Metodologia}

\begin{frame}{Metodologia}

\begin{enumerate}
    \item Revisão da literatura;
    \item Definição do problema;
    \item Construção do modelo;
    \item Implementação computacional;
    \item Execução dos experimentos;
    \item Análise dos resultados.
\end{enumerate}

\end{frame}

% =================================================
\begin{frame}{Modelo matemático}

Considere a variável de decisão:

\[
x_{ij} =
\begin{cases}
1, & \text{se a tarefa } i \text{ for atribuída ao desenvolvedor } j,\\
0, & \text{caso contrário.}
\end{cases}
\]

\vspace{0.3cm}

Exemplo de função objetivo:

\[
\min z =
\sum_{i \in I}
\sum_{j \in J}
c_{ij}x_{ij}
\]

sujeito a:

\[
\sum_{j \in J}x_{ij}=1,
\qquad
\forall i \in I
\]

\end{frame}

% =================================================
\begin{frame}{Fluxo da metodologia}

\centering
\small

\fbox{\strut Dados}

\vspace{0.15cm}

$\downarrow$

\vspace{0.15cm}

\fbox{\strut Modelagem}

\vspace{0.15cm}

$\downarrow$

\vspace{0.15cm}

\fbox{\strut Otimização}

\vspace{0.15cm}

$\downarrow$

\vspace{0.15cm}

\fbox{\strut Experimentos}

\vspace{0.15cm}

$\downarrow$

\vspace{0.15cm}

\fbox{\strut Análise dos resultados}

\end{frame}

% =================================================
% RESULTADOS
% =================================================
\section{Resultados}

\begin{frame}{Resultados}

\begin{table}
\centering

\begin{tabular}{lccc}

\toprule
\textbf{Cenário}
& \textbf{Método A}
& \textbf{Método B}
& \textbf{Método C}
\\

\midrule

Cenário 1 & 10.2 & 8.5 & 7.9 \\
Cenário 2 & 12.1 & 9.3 & 8.8 \\
Cenário 3 & 15.4 & 11.7 & 10.2 \\

\bottomrule

\end{tabular}

\end{table}

\vspace{0.4cm}

\begin{block}{Principal resultado}
Destacar aqui o resultado mais importante encontrado.
\end{block}

\end{frame}

% =================================================
\begin{frame}{Análise dos resultados}

\begin{itemize}
    \item O método A apresentou...
    \item A estratégia B reduziu...
    \item O cenário C demonstrou...
    \item Os resultados sugerem que...
\end{itemize}

\end{frame}

% =================================================
% CONCLUSÕES
% =================================================
\section{Conclusões}

\begin{frame}{Conclusões}

\begin{itemize}
    \item O problema foi modelado utilizando...
    \item Os experimentos mostraram que...
    \item A abordagem proposta permitiu...
    \item Os resultados indicam...
\end{itemize}

\vspace{0.4cm}

\begin{block}{Conclusão principal}
Escrever aqui, em uma ou duas frases, a principal conclusão do trabalho.
\end{block}

\end{frame}

% =================================================
\begin{frame}{Trabalhos futuros}

\begin{itemize}
    \item Avaliar novos conjuntos de dados;
    \item Adicionar novas restrições;
    \item Considerar incerteza nos parâmetros;
    \item Comparar novas técnicas de otimização;
    \item Aplicar o modelo em um ambiente real.
\end{itemize}

\end{frame}

% =================================================
% REFERÊNCIAS
% =================================================
\section{Referências}

\begin{frame}{Referências}

\footnotesize

\begin{thebibliography}{99}

\bibitem{ref1}
SOBRENOME, Nome.
\newblock Título do artigo.
\newblock \textit{Nome do Periódico}, 2025.

\bibitem{ref2}
SOBRENOME, Nome.
\newblock Título do trabalho.
\newblock Universidade, 2024.

\bibitem{ref3}
SOBRENOME, Nome.
\newblock Título do livro.
\newblock Editora, 2023.

\end{thebibliography}

\end{frame}

% =================================================
% FINAL
% =================================================
\begin{frame}

\centering

\vspace{1cm}

{\Huge \textbf{Obrigado!}}

\vspace{1cm}

{\Large Perguntas?}

\vspace{1cm}

\texttt{seu.email@universidade.br}

\end{frame}

\end{document}